\documentclass{article}
\usepackage[a4paper,margin=22mm]{geometry}
\usepackage[T1]{fontenc}
\usepackage{graph-separation}
\usepackage{url}
\usepackage{longtable}
\usepackage[hidelinks]{hyperref}
\begin{document}
\title{Graph separation package}
\author{Qingyuan Zhao\thanks{Statistical Laboratory, University of Cambridge,
qyzhao@statslab.cam.ac.uk.}}
\maketitle
\noindent Version 3.0.0-rc.2: B2 design, draft of 29 October 2026.

Source and local workflow:\\
\url{https://github.com/causal-design/graph-separation-latex}.

\section{Installation}

Download \texttt{graph-separation.sty} and put it beside your \texttt{.tex} file.
Add this line to the preamble:
\begin{verbatim}
\usepackage{graph-separation}
\end{verbatim}
For a shared installation, place the style file under
\texttt{tex/latex/graph-separation/} in your personal TEXMF tree. Locate that tree
with \verb|kpsewhich -var-value=TEXMFHOME|. Refresh the filename database with
\texttt{mktexlsr} if your TeX installation requires it. Do not edit the system's
package directories. A copy beside the document is the simplest option and also
works in Overleaf: upload the style file into the project.

The package requires \texttt{amsmath}, \texttt{amssymb}, \texttt{tikz},
\texttt{adjustbox} and \texttt{centernot}, supplied by standard TeX distributions.

Compile with pdfLaTeX, LuaLaTeX or XeLaTeX; shell escape is not needed. This draft
targets traditional LaTeX mathematics. It has not been qualified for
\texttt{unicode-math}, arbitrary replacement math fonts or automatic bold symbols.
The package draws notation inside equations; it is not a graph-layout package.

\section{Quick start}
Save the following as \texttt{example.tex}, beside the style file, and compile
with \texttt{pdflatex example.tex}:
\begin{verbatim}
\documentclass{article}
\usepackage{graph-separation}
\begin{document}
An inline relation: $V_j\rdpath V_k$.
\[
  V_j\mconnarc V_k \qquad V_j\dconnarc V_k
  \qquad V_j\trek V_k \qquad V_j\notrek V_k.
\]
\end{document}
\end{verbatim}
The display above produces:
\[
  V_j\mconnarc V_k \qquad V_j\dconnarc V_k
  \qquad V_j\trek V_k \qquad V_j\notrek V_k.
\]
Use all commands in math mode. They are mathematical relations and supply their
own relation spacing: do not add manual spaces around a single symbol.

\section{Design}
The selected B2 design is drawn with TikZ, not assembled from AMS glyphs. At a
10\,pt mathematical reference size, all seven squiggly primitives have the same
13.18785\,pt advance and intended ink span, with a 0.45\,pt body stroke. This span
is 90\% of the former 0.515\,cm path length. The wave has two complete triangular
cycles of amplitude 0.8\,pt, with 3.2\,pt straight end clearances. The original
full and half \texttt{Classical TikZ Rightarrow} heads retain a 1\,mm length;
their pen follows TikZ's original head construction.

The wave is identical for every endpoint combination. Arrowheads attach to
separate straight stubs, so adding a head does not shorten or rephase the wave.
All dimensions scale proportionally with the active mathematical style, including
subscripts and sub-subscripts. The centre line sits on the mathematical axis.

\textbf{Symmetry means rotation, not reflection.} The wave, the unheaded symbol,
the full/full symbol and the half/half symbol coincide with themselves under a
180-degree rotation. Right and left variants rotate into each other, as do
half/full and full/half. Half-heads point up at the left end and down at the right.
Superscript labels remain upright: the rotational property describes the
underlying relation glyph, not its \(\mathrm{d}\) or \(\mathrm{t}\) annotation.

Straight primitives retain the published website/course v2.0 drawing, including
its original length and 0.5\,pt stroke. They are intentionally not the same length
as the new squiggles. Existing command names and aliases are retained, but
changed squiggle dimensions can alter line breaks in older documents.

\section{Primitive commands}
Names describe the left endpoint, line kind and right endpoint:
\texttt{no}, \texttt{half} and \texttt{full} name the endpoint marks;
\texttt{straig} and \texttt{squig} distinguish straight and squiggly lines.

\begin{center}
\begin{tabular}{llll}
$\nostraigno$ & \verb|\nostraigno| & $\nosquigno$ & \verb|\nosquigno|\\[5pt]
$\nostraigfull$ & \verb|\nostraigfull| & $\nosquigfull$ & \verb|\nosquigfull|\\[5pt]
$\fullstraigno$ & \verb|\fullstraigno| & $\fullsquigno$ & \verb|\fullsquigno|\\[5pt]
$\halfstraighalf$ & \verb|\halfstraighalf| & $\halfsquighalf$ & \verb|\halfsquighalf|\\[5pt]
$\halfstraigfull$ & \verb|\halfstraigfull| & $\halfsquigfull$ & \verb|\halfsquigfull|\\[5pt]
$\fullstraighalf$ & \verb|\fullstraighalf| & $\fullsquighalf$ & \verb|\fullsquighalf|\\[5pt]
$\fullstraigfull$ & \verb|\fullstraigfull| & $\fullsquigfull$ & \verb|\fullsquigfull|
\end{tabular}
\end{center}
\section{Edges, walks and connections}
\noindent The higher-level commands describe edges, walks, arcs and connections.
An asterisk combines pieces of a walk; the labels \(\mathrm{d}\) and
\(\mathrm{t}\) distinguish d-connection and treks. In particular,
\verb|\mconnarc| has no superscript \(\mathrm{m}\).
\begin{center}
\begin{longtable}{llll}
$\udedge$ & \verb|\udedge| & $\noudedge$ & \verb|\noudedge|\\[5pt]
$\rdedge$ & \verb|\rdedge| & $\nordedge$ & \verb|\nordedge|\\[5pt]
$\ldedge$ & \verb|\ldedge| & $\noldedge$ & \verb|\noldedge|\\[5pt]
$\bdedge$ & \verb|\bdedge| & $\nobdedge$ & \verb|\nobdedge|\\[5pt]
$\rdpath$ & \verb|\rdpath| & $\nordpath$ & \verb|\nordpath|\\[5pt]
$\ldpath$ & \verb|\ldpath| & $\noldpath$ & \verb|\noldpath|\\[5pt]
$\trek$ & \verb|\tconnarc| or \verb|\trek| & $\notrek$ & \verb|\notconnarc| or \verb|\notrek|\\[5pt]
$\mconnarc$ & \verb|\mconnarc| & $\nomconnarc$ & \verb|\nomconnarc|\\[5pt]
$\dconnarc$ & \verb|\dconnarc| & $\nodconnarc$ & \verb|\nodconnarc|\\[5pt]
$\confarc$ & \verb|\confarc| & $\noconfarc$ & \verb|\noconfarc|\\[5pt]
$\samedist$ & \verb|\samedist| & $\confpath$ & \verb|\confpath|\\[5pt]
$\colliderconn$ & \verb|\colliderconn| & $\markovblanket$ & \verb|\markovblanket|\\[5pt]
$\uconn$ & \verb|\uconn| & $\tconn$ & \verb|\tconn|\\[5pt]
$\mconn$ & \verb|\mconn| & $\dconn$ & \verb|\dconn|\\[5pt]
$\muconn$ & \verb|\muconn| & $\deltaconn$ & \verb|\deltaconn|
\end{longtable}
\end{center}
\section{Labels, negation and composition}
Use \verb|\dconnarc| for the half-headed arc with a \(\mathrm{d}\) label,
and \verb|\trek| (also \verb|\tconnarc|) for the full-headed symbol with a
\(\mathrm{t}\) label. \verb|\mconnarc| has no \(\mathrm{m}\) superscript.
Negated variants are explicit commands, for example \verb|\nordpath|,
\verb|\nomconnarc|, \verb|\nodconnarc| and \verb|\notrek|.
\[
 A\nordpath B\quad A\nomconnarc B\quad A\nodconnarc B\quad A\notrek B.
\]
The composite commands retain the original asterisk-separated notation:
\[
 A\mconn B\qquad A\dconn B\qquad A\tconn B.
\]
For example, \verb|$A\mconn B$| renders the half/full piece, an asterisk,
and the full/half piece. The command table documents the available notation;
mathematical definitions and assumptions belong to the accompanying paper or
course. The package itself does not test graph separation.

\section{Smaller styles and compatibility}
No special commands are needed inside subscripts:
\begin{verbatim}
$X_{A\mconnarc B}\qquad X_{Y_{A\trek B}}$
\end{verbatim}
\[
 X_{A\mconnarc B}\qquad X_{Y_{A\trek B}}.
\]
The squiggles scale with the mathematical font's quad rather than using separate
optical designs. Straight primitives keep their legacy fixed-size behaviour.
Labels, negations and composite commands are built from the same primitives.
Do not load another local copy or an old preamble defining these command names:
that will cause duplicate-definition errors. Use \texttt{kpsewhich
graph-separation.sty} to check which installed copy TeX finds; the compile log
should identify version 3.0.0-rc.2.

\section{HTML and reproducible builds}
MathJax cannot execute TikZ directly. The accompanying local Python workflow
renders the native relations through LaTeX and \texttt{dvisvgm}, then creates a
self-contained MathJax header with SVG images and measured widths/baselines.
It keeps the original macro names in editable equations and selects separate
renders for mathematical styles. The package source is the authority; do not
maintain a separately drawn HTML version.
\begin{verbatim}
python3 tools/workflow.py check
python3 tools/workflow.py check-outlines
python3 tools/workflow.py build
python3 tools/workflow.py export-mathjax
\end{verbatim}
The workflow requires Python 3.9 or later; SVG export additionally needs
\texttt{standalone} and \texttt{dvisvgm}. A shared consumer should pin the package
and generator hashes, updating deliberately rather than silently changing an
existing publication. The former AMS draft is superseded. Publishing this source
package does not replace existing website downloads or immutable lecture releases.

\section{Licence}
Copyright 2024, 2026 by Qingyuan Zhao. The package, documentation and local
workflow are released under the MIT License. You may use, modify and redistribute
them, including commercially, subject to retaining the copyright and permission
notice. The software is provided without warranty. See \texttt{LICENSE} for the
full terms. Historical comparison fixtures retain their original notices.
\end{document}
